\documentclass[12pt]{article}

\usepackage[margin=1in]{geometry}
\usepackage{setspace}
\usepackage{amsmath}
\usepackage[style=authoryear,backend=biber]{biblatex}
\addbibresource{lifecycle.bib}
\usepackage{pgf}

\begin{document}

\title{Housing Wealth Over the Life Cycle}
\author{Thomas Davidoff\thanks{Sauder School of Business, University of British
Columbia}}
\maketitle

\begin{abstract}
This chapter summarizes facts and economic theory regarding the accumulation and
disposal of housing debt and equity over consumers' life cycles.  After
reviewing some basic statistics concerning housing wealth in the United States,
I turn to the body of research concerning the way that homeownership affects
households' ability to allocate resources to the dates and states of the world
in which they are most needed. Despite decades of marketing efforts, financial
products meant to undo certain portfolio effects of homeownership, such as
equity sharing and home price index trading, have yet to become popular. A
striking trend in North America is diminished household headship and ownership
among younger people. 

\end{abstract}

\section{Housing Wealth Over the Life Cycle: US Statistics}

Data from surveys of US consumers highlight the importance of housing as
both an asset and a liability. \textcite{freddiemac2024homeowner}, summarizing
U.S. Consumer Expenditure Survey panel data, find that between 1984 and 2023,
rent grew fairly steadily as a share of renters' total expenditures, from
roughly 18\% to 26\%. By contrast, owners expenditures on shelter (including
mortgage interest but not principal repayment) were a relatively constant share
of expenditures, hovering around 15\% over time over that period.\footnote{The
	opportunity cost of equity and anticipated capital gains, part of true
	``user cost'' (see, e.g.  \textcite{Poterba}), are not included in
owners' expenditure shares.}

The greater housing expenditure share for renters than owners found in the CEX is
consistent with the empirical consensus that the wealth elasticity of housing
demand is less than one (that is, the fraction of wealth spent on housing falls
with wealth).\footnote{See, e.g. \textcite{Carliner}.}  In the 2000 IPUMS census
data, I find an elasticity of rent paid with respect to social security income
(intended as a proxy for lifetime income) of just .3 among retirees living
alone who rent housing aged between 65 and 76.  In the 2020-2024 5-year ACS,
that value rises to .47.\footnote{These results are conditional only on
metropolitan area dummies and the year for the 5-year ACS.} Rent relative to
social security income is greater in larger and more expensive metropolitan
areas, consistent with the nearly experimental result of
\textcite{HanushekQuigley} that the price elasticity of demand for housing is
less than one (that is, spending on housing rises as the price of housing
rises). Despite similar incomes, renters in central cities in the 2007 CEX spend
22\% of income on shelter, as opposed to 13\% among renters in rural
areas.\footnote{An intriguing result of \textcite{DavisFOM} is that despite
	these low income and price elasticities, the distribution of
expenditure shares among renters is very close to constant across time and US
metropolitan areas.}

Tabulations of 2004 Survey of Consumer Finances data in Table \ref{tab:assets}
show that for most homeowners of every age group, housing represents a majority
of assets.\footnote{These tabulations were performed by
Todd Sinai and Moises Yi.} For most owners under 50, the value of the home is
greater than net worth. The ratio of home value to assets and net worth is
declining with age and income. These patterns continued to hold in the 2016
SCF, per \textcite{nahb2018homeownership}.

\begin{table}
	\caption{\label{tab:assets} Home value as a fraction of assets (top
	panel) and net worth (bottom panel) for homeowners only, median
	values across age and income groups.  Negative net worth households
	excluded from the Net Worth panel.  Data from the 2007 Survey of
	Consumer Finances, compiled by Todd Sinai and Moises Yi}
	\center{\textbf{Median Home Value as Fraction of Assets}}
	\begin{tabular}{l|llllll}
		\hline
		&                   \multicolumn{5}{c}{Household income}\\ 
		Age of Head      &     $<$40k    & 40-75k   &75-125k  &125-250k    &250+ k
		&Total\\
		\hline\hline
		 under 35& 1.890& 1.848& 1.456& 1.311& 0.576& 1.692\\
		 35-49& 1.469& 1.355& 0.984& 0.761& 0.271& 1.063\\
		 50-64& 0.944& 0.697& 0.572& 0.529& 0.260& 0.634\\
		 65+& 0.752& 0.451& 0.439& 0.265& 0.203& 0.627\\
		\hline
	\end{tabular}
	\\
	\center{\textbf{Median Home Value as Fraction of Net Worth}}
	\\
	\begin{tabular}{l|llllll}
		\hline
		&                   \multicolumn{5}{c}{Household income}\\ 
		Age of Head      &     $<$40k    & 40-75k   &75-125k  &125-250k    &250+ k
		&Total\\
		\hline\hline
		 under 35 &2.188& 1.921& 1.562& 1.311& 0.576& 1.798\\
		35-49 &1.488& 1.367& 0.990& 0.761& 0.271& 1.070\\
		50-64 &0.944& 0.697& 0.572& 0.529& 0.260& 0.637\\
		  65+ &0.754& 0.451& 0.439& 0.265& 0.203& 0.631\\
		Total &0.918& 0.996& 0.829& 0.625& 0.258& 0.849\\
		\hline
	\end{tabular}
\end{table}

Figure \ref{fig:ageMarriageOwnershipIPUMS} plots homeownership and marriage
rates among household heads by age for the 2000 Census 5\% sample and the
2020-2024 ACS based on data drawn from \textcite{IPUMS}.  While overall homeownership
has been relatively constant between the two samples (71\% for 2000 and 69\%
for 2020-2024), we see large declines between 2000 and 2020-2024 at younger
ages (even conditional on household headship) and even more persistent
homeownership among seniors. 

\begin{figure}
	\caption{\label{fig:ageMarriageOwnershipIPUMS} Homeownership and marriage rates by age for all US adults, 2000 Census and 2020-2024 5-year American Community Survey, via \textcite{IPUMS}.}
	\pgfimage[height=8cm]{ageMarriageOwnership.png}
\end{figure}


Both homeownership and marriage rates rise steeply in the 20s and 30s, with
marriage rates levelling off at a younger age. Married couples (with spouse
present) are far likelier to own than rent housing, and far likelier than other
individuals to own (82\% versus 55\% in the 2020-2024 ACS, per \textcite{IPUMS}
calculations). This fact is consistent with the finding of \textcite{VentiWise}
and others that death of a spouse is among the most important triggers of exit
from homeownership among the elderly. \textcite{CoileMilligan} use HRS/AHEAD
data to show that the downward trend in homeownership in old age accelerates
into the 90s, with ill health of individuals joining death of a spouse as a
critical cause. They find an ownership rate among individuals of approximately
50\% at age 90, but approximately 35\% at age 93.

Figure \ref{fig:mortgageStatusByAge} plots the rate of mortgage indebtedness (a
0-1 indicator) for homeowners by age of household-heading homeowners in the
2000 US Census 5\% sample  and the 2020-2024 5-year ACS (both from
\textcite{IPUMS}). In both samples, not surprisingly, the presence of mortgage
debt is very high (over 80\%) among mid-20s homeowners, and declines with age
well into retirement. Somewhat surprisingly in light of recent increases in
home prices, mortgage debt has grown less common among younger homewners and more
common among older homeowners. The overall incidence among owners has fallen
from 66\% in 2000 to 59\% in the 2020-2024 ACS. 

\begin{figure}
	\caption{\label{fig:mortgageStatusByAge} Presence of mortgage debt by age for household-heading homeowners, 2000 US Census 5\% sample and 2020-2024 5-year ACS, via \textcite{IPUMS}.}
	\pgfimage[height=8cm]{mortgageStatusByAge.png}
\end{figure}
 
Making use of the patterns depicted in Figure \ref{fig:ageMarriageOwnershipIPUMS},
\textcite{FisherGervais} show that homeownership rates declined sharply among
younger households between 1980 and 2007, attributing this in large part to
increasing age of first marriage over time. Working in the other direction,
\textcite{Kearney} finds that rising home prices are a cause of reduced
fertility (a major trend globally). \textcite{StatsCan2024} shows that in
Canada, where home prices rose dramatically over past decades, while
Millennials represent a large population bulge at prime parenting ages in 2024
relative to their predecssors in 2004, the number of children has changed
relatively little over the same period. Whether the receding of the Millenial
bulge will lead to lower home prices and hence increased marriage and fertility
rates remains to be seen.

Gifts from parents to children have long been important sources of the
downpayments lender typically require of homeowners (see, e.g.
\textcite{englehardtMayer}). As home prices have risen in Canada, the fraction
of first-time buyers purchasing with gifts from parents has risen from 20\% in
2015 to 31\% in 2024, per \textcite{Tal}. Over the same period the average gift
rose from under \$60,000 to over \$100,000. Home equity loans are a natural source of these gifts, suggesting a channel for the rising US indebtedness among seniors shown in Figure \ref{fig:mortgageStatusByAge}.

Increasing overall homeownership can be attributed largely to an aging
population and increasing homeownership among the elderly. These patterns have
persisted. Per 2000 Census data from \textcite{IPUMS}, among respondents aged
75 and over, in 2000, 69.8\% were household heads or spouses of heads and owned
their own residence. In the 2024 ACS, it was 72.9\%. Among Americans 24-35, in
2000, 46\% were household heads or spouses of and owned their home. In 2024,
that figure had dropped to 34.9\%. This reflects both a reduction in household
headship rates and the conditional probability of ownership. While rising
housing costs are a global phenomenon, \textcite{OECDLaunch} shows that between
2007 and 2019, the share of individuals aged 20-29 living with parents was
roughly constant overall for OECD countries, but rose dramatically along with
home prices in the U.S. and Canada.

Among US households, the median and modal number of residential properties owned
is one. \textcite{AizcorbeKennickellMoore} find that 67\% of households owned
housing in 2001, but only 12.8\% owned a second home or rental housing.  

The unconditional mean income of owners in the 2020-2024 American Community
Survey was \$101,988, almost double the mean income of renters (\$57,843). In
the 2000 Census (IPUMS microdata), the median homeowner household income was
\$49,600, as opposed to the median renter income of \$26,800. This difference
in incomes by ownership is consistent with a number of prominent explanations,
such as: the US tax code is increasingly favorable to homeownership as income
and wealth rise, the normality and age profile of demand for detached homes
combined with a natural tendency for detached units to be owner
occupied,\footnote{See, e.g. \textcite{GlaeserShapiro}.} and a greater prevalence
of credit constraints among those with lower income and wealth. 

Consistent with declining tax benefits to homeownership in the U.S. due to
elimination of state and local tax deductions as well as reduced mortgage
deductibility benefits due to falling rates, rising standard deductions, and
caps on mortgage interest deductions, the relationship between homeownership
and income weakened slightly over time, with the regression coefficient on log
income (with an indicator for homeownership on the left hand side) falling from
.099 in 2000 Census data (taken from \textcite{IPUMS}) to .084 in 2020-2024 ACS
data (taken from \textcite{IPUMS}). Not surprisingly given the positive
geographic cross-sectional relationships among education, home prices, income,
and regulations (\textcite{davidoffCFR}), these coefficients rise to .113 for 2000 and
.096 for 2020-2024, in the presence of controls for metropolitan area fixed
effects.\footnote{The full regression specification is a linear probability OLS
regression of log income on age dummies and marital status dummies, with the
sample selected on household heads ages 25 and older.}


\section{Housing wealth in the life cycle: theoretical problems}

The institution of homeownership presents a puzzle.  Why do most households own
exactly 100\% of a single home, rather than some smaller share of a diversified
bundle of homes? A variety of costs come with full ownership of a particular
home, some of which are described below.  

\subsection{Imposition of mobility costs}

For a given household, choosing to own rather than rent housing makes moving
costlier. Finding a new occupant for a home requires maintenance and search
effort independent of tenure. However, brokerage and transfer fees are
significantly larger for owner units, and more marketing effort is presumably
required when the next occupant is likelier to be an owner than a renter, and
hence expect a longer stay in the home. Purchasing a unit suited for owner
occupancy will thus render mobility costly, even if lifetime turnover costs are
lower for owner units.

Between the effects of \emph{ex ante} selection and  \emph{ex post} costs, it is
clear empirically that owners are less mobile than renters.
\textcite{SchachterKuenzi} show that in the 1996 Current Population Survey, the
distribution of length of time in current residence of owners first order
dominates that of renters.  The same is true in the 2000 US census.
\textcite{ChettySzeidlra} find that in response to unemployment shocks, renters are
more likely to move than owners, both in levels and relative to baseline
mobility.  Also, homeowners decrease non-housing consumption more in response to
unemployment than do renters.  To the extent that housing and other goods are
poor substitutes, ownership may thus exacerbate the expected utility costs of
the volatility of wealth.  In the 2000 US Census and 2020-2024 5-year ACS (per
\textcite{IPUMS}), conditional on income, the number of bedrooms in
respondents' units declines significantly more rapidly with age among renters
than owners, consistent with a greater cost of downsizing for elderly owners.

An additional factor rendering homeowners immobile at times in the U.S. is the
prevalence of fixed rate mortgage loans that can be refinanced without cost
downward. When rates rise, numerous studies (e.g. \textcite{Rothstein}) show
that owners' mobility is reduced through a lock-in effect.

Whether the existence of demand for changing units makes owning or renting
housing relatively advantageous is an open question.  The answer will depend in
large part on the relative magnitudes of: (a) moral hazard costs of rapid
turnover and under-maintenance that arise because tenants do not pay on the
margin for landlords' leasing and marketing costs upon lease termination, versus
(b) the expected extra financial costs of moving imposed by owner units.  The
cost (a) is a common rationalization of the institution of homeownership through
the mechanism of depreciation and maintenance.\footnote{See, e.g.
\textcite{Shelton}.}  The magnitude of cost (b) will depend in part on households
ability to substitute consumption between housing and other goods and between
current housing and future expenditures.  To the extent that homeowners purchase
homes that are particularly well-suited to their idiosyncratic needs or
attachment to the home grows over time in a beneficial way (for example, through
customization) cost (b) will be smaller.

\subsection{Excessive investment in particular homes}

\textcite{HendersonIoannides} observe that there is a welfare cost to homeownership
if investment demand for a consumer's optimally chosen home is less than
consumption demand for that home (in the opposite case, one can invest in
housing assets beyond one's own home). There are two obvious ways in which
realistic market imperfections will render investment demand relatively weak.
First, homeowners may need to hold home equity greater than the amount of total
savings they would otherwise choose. Second, homeowners may allocate too much
portfolio weight to housing in general and the price risk of their own home in
in particular.

\subsubsection{Too much saving early and late in life}

Income for an individual or couple is typically hump shaped over the life cycle.
If a flat trajectory of expenditures is desirable, then transfers from middle
age to earlier and later years are also desirable.  As emphasized by
\textcite{ArtleVaraiya}, homeownership typically involves saving for a down payment
in youth, and retaining home equity through much of retirement.  Both of these
characteristics of homeownership interfere with transferring resources from
middle age to the early and late years of low earnings.

Both casual and formal empirical analysis demonstrate that down payment
constraints are undesirable from the perspective of younger households.  A
first piece of evidence is the relative popularity of very high loan to value
mortgages among younger households (shown in recent Canadian loan data by
\textcite{AlAboud2025}). Teaser loans, popular during the runup to the Great
Financial Crisis featured high payments in the later years of the loan's life,
suggesting a willingness to pay for the convenience of avoiding a large down
payment. That housing prices increased sharply in response to expansion of
mortgage credit in that period, and crashed thereafter, also suggests that
credit constraints are binding, as in the model of \textcite{OrtaloMagneRady}.  

Using microdata covering years in which high loan to value ratios were
unavailable or extremely costly, \textcite{Engelhardt96} shows that consumers spend
less money on food in the years leading up to a first purchase of housing than
they do in subsequent years. This again suggests that there are significant
liquidity constraints associated with down payments.

The elderly retain a strikingly high quantity of home equity, a fact that has
been documented by numerous authors such as \textcite{VentiWise}. As shown in
figure \ref{fig:homeownershipAge_allMetros}, home ownership rates remain high
very late into life, and mortgage debt is decreasing, not increasing in age.
As reported in \textcite{Davidoffltc}, in the 2004 wave of the HRS/AHEAD panel,
approximately 12\% of retired homeowners over 62 owe any mortgage debt, and
among these, the median debt to value is just 33\%.  Between 1998 and 2004,
despite very large increases in home values, only approximately 10\% of retired
homeowners increased mortgage debt.  Almost all high wealth households have
substantial home equity.  In the top total wealth quintile of the retired and
over 62 sample, 84\% held home equity greater than \$100,000.  To the extent
that retirees wish to spend down their housing wealth before death, it thus
appears that a combination of the institution of homeownership and supply
problems in the markets for residential sale-leasebacks and mortgages for the
elderly stand in the way. Following a peak after the Great Financial Crisis,
reverse mortgage originations have droppped below termination rates.\footnote{See \textcite{NRMLA2026}, home equity release products in the UK appear to have grown more than elsewhere per \textcite{MayerMoulton}.}

That there is a desire to spend down financial resources before death is not at
all obvious, as demonstrated by ambiguous findings on the overall savings rate
of the elderly, e.g.  \textcite{Hurd90}, \textcite{Horioka} and \textcite{Weil}.  Even
without a desire to bequeath wealth to heirs, a low level of consumption in
retirement may be justified by a need for precautionary savings. Precautionary
savings might arise for a variety of reasons, such as a stochastic need for
long-term care and uncertain longevity, combined with imperfect insurance
markets.  Indeed, \textcite{Davidoffpuzzle} shows that the welfare benefits of
each one of home equity extraction products, annuities, and long-term care
insurance depend on the presence or absence of the others.\footnote{For a range of views on home equity crowd-out of long-term care insurance on empirical, survey, and theoretical grounds, see \textcite{Achou}, textcite{Hanewald}, and \textcite{Hirth}.}

A critical input to the retiree's financial planning problem is the question of
how costly in financial and utility terms it would be to move out of one's home.
If it is easy to move out of the home, then a reverse mortgage is presumably
less desirable, and the need to maintain a buffer stock of other assets is
lessened.  Survey evidence in \textcite{AARP} suggests that moves are quite costly.
Engelhardt shows that random increases in social security income and support
for home health care provision increase homeownership. \textcite{CoileMilligan}
show that ill health is a crucial determinant of exit from homeownership.  A
reasonable framework for analyzing homeownership in old age would thus pose
a trade-off between the liquidity and potential financial gain to selling a home
versus the health-dependent utility from remaining in place.  Calibration of
such a model for a given household to provide normative advice on savings and
optimal move date would stretch economists' tools.

\subsubsection{Excessive portfolio weight on a particular home}

The possibility that homeownership might lead to excessive portfolio weight on
housing is been the subject of growing academic, industry, and government interest over
recent decades (see, e.g.  \textcite{FlavinYamashita}, \textcite{Brueckner},
\textcite{Bula}).  Setting aside moral hazard considerations described above,
it is hard to believe that massive exposure to the considerable idiosyncratic
risk embodied in a particular house would be part of an optimal allocation.
Establishing an optimal portfolio share for a long position in some market's
housing prices in general, however, is extremely difficult. I outline some of
the difficulties in determining this optimal share below.

To determine the optimal portfolio exposure to home prices, one must be able to
determine whether an investor is better or worse off if home prices rise.  The
sign of the welfare effect of a price increase depends in part on the extent to
which elderly homeowners plan to use home equity to finance retirement.  If
homeowners simply purchased a home when young, and then used the proceeds to
finance retirement, and if rents were uncorrelated with the price of a home,
then increasing price variability would reduce portfolio demand for housing, and
increasing the size of the home would make the portfolio riskier.  

At the other extreme, if housing demand is constant forever and prices are
perfectly correlated with rents, then owning a home is a perfect hedge against
expenditure risk, and increases in price variability would increase the value of
homeownership relative to renting.  Given the ability to substitute between
housing and other consumption in an otherwise steady state world, a familiar
result is that a one-time change to housing prices must be weakly welfare
increasing for a homeowner.

Not just hedging demand, but also real options associated with timing of moves,
investment, default, and mortgage refinancing may make price volatility more
tolerable to owners at any stage of the life cycle.  In the recent boom, it
seems highly likely that demand and prices were increased by owners and
investors in housing who viewed high loan-to-value loans as essentially call
options on volatile local prices.

The fact that many homeowners die without cashing out their home equity may be
seen as evidence that capital gains risk is of limited importance.  However, if
housing is held as a buffer for unlikely events, price variation may impose
considerable costs as the states in which the home is sold and prices are low
may feature particularly high marginal utility.  

Just as volatility of prices may increase or decrease the benefits of purchasing
more or less housing, volatility of rents has ambiguous effects on the welfare
of renters.  If housing and other goods are highly substitutable within and
between periods, then variability of rents may increase the welfare of renters
(this is obvious in the case where housing and other goods are perfect
substitutes).  In this case, there may be no hedging demand for homeownership.
That demand for housing is sufficiently flexible to rule out a motivation to
hedge future expenditures runs counter to intuition, and seems unlikely given
price elasticity estimates below one.

Among older households, \textcite{Skinner} observes that it is difficult to
distinguish the relative importance of bequest motives, precautionary
considerations, and optimal risk averse exercise of an option to sell in the
retention of home equity in old age.  Indeed the analysis and caveats of
\textcite{EvansHendersonHobson} suggest that just the relationship between optimal
date of sale and price volatility is almost impossible to sign.  Some evidence
that capital gains risk matters, supportive of a precautionary motive, is
provided by \textcite{BanksBlundellOldfieldSmith} and \textcite{SinaiSouleles}.  Both
find that older households tend to sell out of high price volatility
environments in old age.  The authors' favored interpretation of this evidence
is that hedging considerations dominate in youth, so that volatility may boost
portfolio demand for housing, but capital gains risk dominates in old age so
that volatility decreases portfolio demand for housing late in the life cycle.

An easy test of the notion that homeowners are exposed to too much average price
risk comes from the fact that consumers who recognize that they have excessive
exposure to housing in their portfolios would seek to avoid exposing themselves
to further home price risk with their other investment choices.  To my
knowledge, there is no study of the cross sectional relationship between (a) the
correlation of assets with real estate returns and (b) their relative prices.
Such a study would face the challenge that real estate has historically had a
weak correlation with financial asset prices, but crashed with the stock market
in the great depression and in the recent downturn.  

Active investors may or may not demand compensation for exposure to real estate
price risk, but these investors are unusual.  Findings from the Survey of
Consumer finances reported by \textcite{AizcorbeKennickellMoore} show that many
homeowners are not active investors in other markets.  As of 2001, approximately
50\% of households owned no equities, and many of those who did were passive
investors in managed funds.  Even unpriced exposure to systematic risk in
housing markets does not imply that homeowners have an optimal share of their
portfolio in their own particular home.  Home prices are imperfectly correlated
across markets, so exposure to national or international housing markets is
different from exposure to individual market risk.  Moreover, home price
appreciation varies considerably from neighborhood to neighborhood within
markets, and even from home to home within neighborhoods.  

\section{Minimizing Investment Constraints}
A variety of products have been put forward that may help young and old
homeowners reduce home equity and to help all owners reduce exposure to
systematic and idiosyncratic price risk.  An important open question is whether
these products have proved unpopular because contracting issues pushes prices
away from actuarially fair, or instead because high savings and large housing
shares in the portfolio are optimal even with fair pricing.

The obvious way to reduce savings while owning a home is to take on a mortgage.
The cost of increasing loan-to-value ratios has varied considerably over time.
Recent experience demonstrates that there may be no interest rate at which 100\%
loan-to-value is profitable to lenders.  From lenders' perspective, both
refinancing and default are crucial to pricing and profitability, but the
optimal exercise of these options for borrowers are difficult to solve even for
economists.  The recent implosion of the mortgage industry and microeconomic
evidence makes clear that we do not have a satisfactory model of default and
prepayment, even conditional on complete histories of prices and interest
rates.\footnote{\textcite{DengQuigleyVanOrder} emphasize the heterogeneity of
borrower behavior across households and environments.}  Credit availability to
younger homeowners hoping to minimize savings while retaining the hedging,
contracting, or psychological benefits of homeownership is thus likely to remain
cyclical.

From the discussion above, it is not clear that the elderly wish to reduce
savings in the form of home equity or otherwise.  The obvious products that
would allow the elderly to extract equity without moving out of their homes are
sale-leasebacks and reverse mortgages.  The latter are somewhat more popular
than the former, but even demand for reverse mortgages has been weak, as noted
above. Effective interest rates for reverse mortgages are quite high, with
origination fees in the U.S. currently up to \$6,000 and historically higher.
Reverse mortgages are likely to feature both adverse selection (witness Mme
Jean Calment who lived a record 122 years and had the French version of a
reverse mortgage) and moral hazard both on the timing of sales (once a property
is worth less than the loan amount, the mortgagor has no incentive to move
because the loans are non-recourse) and on maintenance,\footnote{see, e.g.
\textcite{MiceliSirmans}} it is not obvious that sufficiently high loan-to-value
ratios to trigger mass demand can be attained at profitable and marketable rates.

\textcite{Shiller} and \textcite{CaplinChanFreemanTracy} have proposed mechanisms
whereby homeowners can reduce exposure to variation in home prices.  The Chicago
Mercantile Exchange now allows investors to take long or short positions on
changes in regional average home price indexes, calculated based on a repeated
sale methodology.  Several companies have devised mechanisms whereby outside
investors can be given equity, rather than debt, in the home.  For the latter
type of contract, default risk should be lower than with debt contracts, but
investment and market timing of sale may lead to problems of moral hazard.
Trades on the Case-Shiller CME index avoid homeowner agency problems but do not
provide insurance against individual home price risk.  Given the complicated
effects of home price changes on utility described above, it is not entirely
surprising that there has been very limited demand for either type of product.

Summarizing, it is difficult to provide households with easily understood
normative advice on major life cycle housing decisions, and particularly
difficult to characterize the optimal portfolio share for housing. The
conventions of hump shaped life cycle ownership rates, high average
homeownership rates enduring late into retirement, mortgage financing, and
large unhedged positions in individual homes will be with us for the
foreseeable future.

\printbibliography

\end{document}

